\begin{resumo}[\protect\bfseries Resumo]

    Neste relato, descrevo minha trajetória na \ac{ages}, da AGES I até a AGES IV,
    com o objetivo de mostrar minha evolução técnica e humana em projetos de
    software reais. Na AGES I, atuei no projeto Informativo para Imigrantes e
    Refugiados, contribuindo com o desenvolvimento e testes de funcionalidades de um
    aplicativo mobile para organizar serviços e informações úteis para esse público.
    Na AGES II, partcipei do projeto Let Me Trial, onde participei da modelagem
    e do refinamento do banco de dados, além do desenvolvimento para apoiar a
    pré-triagem de pacientes em ensaios clínicos.

    Na AGES III, meu foco foi a arquitetura e integração no projeto
    Se Doce Fosse, ajudando a conectar decisões de banco, interface e backend para
    manter consistência de entrega. Já na AGES IV, trabalhei com uma postura mais
    integrada de coordenação do backlog e alinhamento com stakeholders, no projeto
    Uma Porto Alegre Alemã, consolidando habilidades de comunicação técnica e
    organização do trabalho do time.

    Ao longo dessas etapas, percebi que para se tornar um Engenheiro de Software qualificado não posso depender
    apenas de escrever código: devo compreender o contexto, saber priorizar,
    negociar decisões e aprender com os bloqueios do time.

  \bigbreak\
  \textbf{PALAVRAS CHAVES:}
  AGES, Engenharia de Software, Desenvolvimento Mobile, Desenvolvimento Web,
  Arquitetura de Software, Aprendizagem, Gestão de Equipe.
\end{resumo}
